\chapter{DL-019 - Theme Rule Selection And Inspection}\label{record-184}

\noindent{\footnotesize\color{BookMuted}\path{docs/design/DL-019-theme-rule-selection-and-inspection.md}\par}\smallskip

\begin{RecordMeta}
id: DL-019
title: Theme Rule Selection And Inspection
status: active
lane: cool-ideas
priority: medium
github_issue: 444
legend: DL

\end{RecordMeta}

Legend: \href{https://github.com/flyingrobots/bijou/blob/1853ad902f6bf219326409e7f743a5a1d6517dbc/docs/legends/DL-design-language.md}{DL - Design Language}

\section{Decision Summary}\label{record-184-decision-summary}

Bijou should adopt the strongest ideas from reactive design-token systems without importing an incompatible dependency or replacing the existing runtime theme contract.

The runtime truth remains:

\begin{CodeBlock}
Theme
  -> ResolvedTheme
    -> TokenGraph
      -> ThemeAccessors
        -> terminal cells, lower modes, and DTCG interop
\end{CodeBlock}

This cycle adds a Bijou-native rule layer for theme authoring and inspection:

\begin{itemize}
\tightlist
\item
  selector rules such as best contrast, minimum contrast, most vivid, least vivid, closest color, and ordered choice
\item
  explicit candidate exclusions for reserved roles such as success and error
\item
  structured inspection facts explaining why a rule chose a token
\item
  dependency facts that Theme Inspector and tests can consume later
\item
  compiled output that remains a normal Bijou \texttt{T\allowbreak{}h\allowbreak{}e\allowbreak{}m\allowbreak{}e} or token graph entry
\end{itemize}

No \texttt{d\allowbreak{}e\allowbreak{}s\allowbreak{}i\allowbreak{}g\allowbreak{}n\allowbreak{}-\allowbreak{}b\allowbreak{}o\allowbreak{}o\allowbreak{}k} runtime, dev, peer, optional, or generated dependency is added. The external package is useful prior art, but its license and general CSS surface do not belong in Bijou's Apache-2.0 zero-dependency core.

\section{Sponsored Human}\label{record-184-sponsored-human}

A Bijou theme author wants to tune primitive palettes while keeping semantic roles stable. They should be able to say ``choose readable text for this surface'' or ``choose a vivid accent that is not success or error'' instead of manually maintaining every derived hex value.

\section{Sponsored Agent}\label{record-184-sponsored-agent}

An agent auditing theme behavior wants deterministic facts for rule inputs, candidate rejection, selected values, dependencies, and contrast proof so it can explain a theme without sampling screenshots or guessing from final RGB values.

\section{Hill}\label{record-184-hill}

A Bijou theme author can define semantic theme roles through pure selector rules, replay the product questions below with tests, compile the rules into the existing theme/token graph contract, and inspect which candidates were considered, rejected, and selected.

\section{Playback Questions}\label{record-184-playback-questions}

\begin{itemize}
\tightlist
\item
  Can a semantic foreground choose a readable color from a primitive palette without hard-coding the final hex value?
\item
  Can a vivid accent selector exclude reserved success and error roles while still proving contrast against its surface?
\item
  Can \texttt{T\allowbreak{}o\allowbreak{}k\allowbreak{}e\allowbreak{}n\allowbreak{}G\allowbreak{}r\allowbreak{}a\allowbreak{}p\allowbreak{}h\allowbreak{}.\allowbreak{}i\allowbreak{}n\allowbreak{}s\allowbreak{}p\allowbreak{}e\allowbreak{}c\allowbreak{}t\allowbreak{}(\allowbreak{})} explain every candidate read by a rule, including selected, eligible, excluded, invalid, and contrast-failed candidates?
\item
  Can malformed rule inputs fail with the missing target, scope, candidate, or exclusion path instead of silently choosing a surprising color?
\item
  Can the rule layer preserve existing \texttt{T\allowbreak{}o\allowbreak{}k\allowbreak{}e\allowbreak{}n\allowbreak{}G\allowbreak{}r\allowbreak{}a\allowbreak{}p\allowbreak{}h\allowbreak{}.\allowbreak{}g\allowbreak{}e\allowbreak{}t\allowbreak{}(\allowbreak{})} behavior for callers that only need a resolved \texttt{T\allowbreak{}o\allowbreak{}k\allowbreak{}e\allowbreak{}n\allowbreak{}V\allowbreak{}a\allowbreak{}l\allowbreak{}u\allowbreak{}e}?
\end{itemize}

\section{Current Truth}\label{record-184-current-truth}

Bijou already has:

\begin{itemize}
\tightlist
\item
  a public \texttt{T\allowbreak{}h\allowbreak{}e\allowbreak{}m\allowbreak{}e} contract with \texttt{s\allowbreak{}e\allowbreak{}m\allowbreak{}a\allowbreak{}n\allowbreak{}t\allowbreak{}i\allowbreak{}c}, \texttt{s\allowbreak{}u\allowbreak{}r\allowbreak{}f\allowbreak{}a\allowbreak{}c\allowbreak{}e}, \texttt{b\allowbreak{}o\allowbreak{}r\allowbreak{}d\allowbreak{}e\allowbreak{}r}, \texttt{u\allowbreak{}i}, \texttt{s\allowbreak{}t\allowbreak{}a\allowbreak{}t\allowbreak{}u\allowbreak{}s}, and \texttt{g\allowbreak{}r\allowbreak{}a\allowbreak{}d\allowbreak{}i\allowbreak{}e\allowbreak{}n\allowbreak{}t} token families
\item
  a \texttt{T\allowbreak{}o\allowbreak{}k\allowbreak{}e\allowbreak{}n\allowbreak{}G\allowbreak{}r\allowbreak{}a\allowbreak{}p\allowbreak{}h} with references, light/dark definitions, transforms, cache invalidation, and subscribers
\item
  \texttt{d\allowbreak{}e\allowbreak{}f\allowbreak{}i\allowbreak{}n\allowbreak{}e\allowbreak{}T\allowbreak{}h\allowbreak{}e\allowbreak{}m\allowbreak{}e\allowbreak{}(\allowbreak{})} for mode-aware token themes
\item
  \texttt{f\allowbreak{}r\allowbreak{}o\allowbreak{}m\allowbreak{}D\allowbreak{}T\allowbreak{}C\allowbreak{}G\allowbreak{}(\allowbreak{})} and \texttt{t\allowbreak{}o\allowbreak{}D\allowbreak{}T\allowbreak{}C\allowbreak{}G\allowbreak{}(\allowbreak{})} for token document interop
\item
  \texttt{d\allowbreak{}o\allowbreak{}c\allowbreak{}t\allowbreak{}o\allowbreak{}r\allowbreak{}T\allowbreak{}h\allowbreak{}e\allowbreak{}m\allowbreak{}e\allowbreak{}(\allowbreak{})} and \texttt{d\allowbreak{}e\allowbreak{}f\allowbreak{}i\allowbreak{}n\allowbreak{}e\allowbreak{}T\allowbreak{}h\allowbreak{}e\allowbreak{}m\allowbreak{}e\allowbreak{}S\allowbreak{}a\allowbreak{}f\allowbreak{}e\allowbreak{}P\allowbreak{}a\allowbreak{}i\allowbreak{}r\allowbreak{}s\allowbreak{}(\allowbreak{})} for concrete contrast checks
\item
  DOGFOOD Theme Lab and Theme Inspector surfaces for concrete token facts
\end{itemize}

The missing piece is rule provenance. Today a derived token can be encoded as a manual hex, a direct reference, or a transform. It cannot yet say:

\begin{itemize}
\tightlist
\item
  choose the highest-contrast color from this palette
\item
  choose the first color that reaches a contrast floor
\item
  choose the most vivid color that is still readable
\item
  exclude reserved role tokens from an accent search
\item
  explain which candidates failed and why
\end{itemize}

\section{Problem}\label{record-184-problem}

Manual semantic token maintenance creates drift:

\begin{itemize}
\tightlist
\item
  a primitive palette can change while semantic foregrounds stay stale
\item
  a vivid accent may accidentally reuse a destructive or success role
\item
  contrast proof happens after values are chosen instead of during selection
\item
  Theme Inspector can show final values but not why those values won
\item
  tests can assert concrete hex values but not the decision rule behind them
\end{itemize}

The answer is not to replace Bijou's runtime theme model. The runtime model already encodes terminal-specific facts: \texttt{N\allowbreak{}O\allowbreak{}\_\allowbreak{}C\allowbreak{}O\allowbreak{}L\allowbreak{}O\allowbreak{}R}, foreground/background tokens, text modifiers, RGB caches, ANSI downsampling, lower-mode behavior, and safe-pair checks.

\section{Scope}\label{record-184-scope}

This cycle adds pure core support for rule-based color selection and inspection.

Included:

\begin{itemize}
\tightlist
\item
  color candidate scopes represented as token paths or explicit color values
\item
  selector functions for contrast, vividness, closeness, ordered position, and minimum contrast
\item
  explicit \texttt{n\allowbreak{}o\allowbreak{}t} exclusions by token path
\item
  selector result facts with candidate pass/fail reasons
\item
  dependency extraction for references and candidate scopes
\item
  \texttt{T\allowbreak{}o\allowbreak{}k\allowbreak{}e\allowbreak{}n\allowbreak{}G\allowbreak{}r\allowbreak{}a\allowbreak{}p\allowbreak{}h} support for rule definitions if that can land without breaking existing callers
\item
  focused tests for selection, exclusion, dependency, and cycle behavior
\item
  design-system docs and changelog notes
\end{itemize}

Out of scope:

\begin{itemize}
\tightlist
\item
  no external \texttt{d\allowbreak{}e\allowbreak{}s\allowbreak{}i\allowbreak{}g\allowbreak{}n\allowbreak{}-\allowbreak{}b\allowbreak{}o\allowbreak{}o\allowbreak{}k} package integration
\item
  no CSS renderer
\item
  no typography, spacing, motion, or duration token categories
\item
  no Figma adapter
\item
  no Theme Lab redesign in this slice
\item
  no element-level renderer provenance
\item
  no app-frame theme-mode UX changes
\end{itemize}

\section{API Direction}\label{record-184-api-direction}

The first public API should be small and structural. Exact names can move in implementation, but the shape should stay close to this:

\begin{CodeBlock}
const graph = createTokenGraph({
  palette: {
    ink: "#0b1020",
    paper: "#f8fafc",
    cyan: "#38bdf8",
    green: "#22c55e",
    red: "#ef4444",
  },
  surface: {
    primary: { fg: "#f8fafc", bg: { ref: "palette.ink" } },
  },
  semantic: {
    primary: bestContrastWith(
      { ref: "surface.primary.bg" }, // token background slot
      scope("palette"),
    ),
    accent: mostVivid(scope("palette"), {
      against: { ref: "surface.primary.bg" },
      minContrast: 4.5,
      not: ["palette.green", "palette.red"],
    }),
  },
});
\end{CodeBlock}

If the fluent helper API is too large for the first slice, the serialized graph definition may use plain discriminated data instead:

\begin{CodeBlock}
semantic: {
  primary: {
    rule: "best-contrast-with",
    target: { ref: "surface.primary.bg" },
    candidates: { kind: "scope", path: "palette" },
  },
}
\end{CodeBlock}

Slot-path references ending in \texttt{.\allowbreak{}b\allowbreak{}g} are part of this cycle's contract. A reference such as \texttt{\{\allowbreak{} \allowbreak{}r\allowbreak{}e\allowbreak{}f\allowbreak{}:\allowbreak{} \allowbreak{}"\allowbreak{}s\allowbreak{}u\allowbreak{}r\allowbreak{}f\allowbreak{}a\allowbreak{}c\allowbreak{}e\allowbreak{}.\allowbreak{}p\allowbreak{}r\allowbreak{}i\allowbreak{}m\allowbreak{}a\allowbreak{}r\allowbreak{}y\allowbreak{}.\allowbreak{}b\allowbreak{}g\allowbreak{}"\allowbreak{} \allowbreak{}\}} resolves the background color of the \texttt{s\allowbreak{}u\allowbreak{}r\allowbreak{}f\allowbreak{}a\allowbreak{}c\allowbreak{}e\allowbreak{}.\allowbreak{}p\allowbreak{}r\allowbreak{}i\allowbreak{}m\allowbreak{}a\allowbreak{}r\allowbreak{}y} token while \texttt{\{\allowbreak{} \allowbreak{}r\allowbreak{}e\allowbreak{}f\allowbreak{}:\allowbreak{} \allowbreak{}"\allowbreak{}s\allowbreak{}u\allowbreak{}r\allowbreak{}f\allowbreak{}a\allowbreak{}c\allowbreak{}e\allowbreak{}.\allowbreak{}p\allowbreak{}r\allowbreak{}i\allowbreak{}m\allowbreak{}a\allowbreak{}r\allowbreak{}y\allowbreak{}"\allowbreak{} \allowbreak{}\}} keeps the existing foreground-reference behavior.

The graph should expose inspection facts separately from resolved values:

\begin{CodeBlock}
const fact = graph.inspect("semantic.accent", "dark");
\end{CodeBlock}

The returned fact should be structured enough for tests and future Theme Inspector rows:

\begin{CodeBlock}
semantic.accent
  rule: most-vivid
  selected: palette.cyan (#38bdf8)
  dependencies: surface.primary.bg, palette.ink, palette.paper, palette.cyan,
    palette.green, palette.red
  rejected:
    palette.paper - eligible, not selected
    palette.green - excluded
    palette.red - excluded
    palette.ink - contrast below 4.5
\end{CodeBlock}

\section{Selector Vocabulary}\label{record-184-selector-vocabulary}

The initial selector vocabulary should cover the high-value color decisions. Selectors that compare against a target put the target first. Selectors that rank only a candidate collection put the collection first.

\texttt{b\allowbreak{}e\allowbreak{}s\allowbreak{}t\allowbreak{}C\allowbreak{}o\allowbreak{}n\allowbreak{}t\allowbreak{}r\allowbreak{}a\allowbreak{}s\allowbreak{}t\allowbreak{}W\allowbreak{}i\allowbreak{}t\allowbreak{}h\allowbreak{}(\allowbreak{}t\allowbreak{}a\allowbreak{}r\allowbreak{}g\allowbreak{}e\allowbreak{}t\allowbreak{},\allowbreak{} \allowbreak{}c\allowbreak{}a\allowbreak{}n\allowbreak{}d\allowbreak{}i\allowbreak{}d\allowbreak{}a\allowbreak{}t\allowbreak{}e\allowbreak{}s\allowbreak{})} chooses the candidate with the highest contrast ratio against the target. Equal scores keep candidate order.

\texttt{m\allowbreak{}i\allowbreak{}n\allowbreak{}C\allowbreak{}o\allowbreak{}n\allowbreak{}t\allowbreak{}r\allowbreak{}a\allowbreak{}s\allowbreak{}t\allowbreak{}W\allowbreak{}i\allowbreak{}t\allowbreak{}h\allowbreak{}(\allowbreak{}t\allowbreak{}a\allowbreak{}r\allowbreak{}g\allowbreak{}e\allowbreak{}t\allowbreak{},\allowbreak{} \allowbreak{}c\allowbreak{}a\allowbreak{}n\allowbreak{}d\allowbreak{}i\allowbreak{}d\allowbreak{}a\allowbreak{}t\allowbreak{}e\allowbreak{}s\allowbreak{},\allowbreak{} \allowbreak{}\{\allowbreak{} \allowbreak{}r\allowbreak{}a\allowbreak{}t\allowbreak{}i\allowbreak{}o\allowbreak{} \allowbreak{}\}\allowbreak{})} chooses the first candidate that clears the requested contrast ratio, using deterministic candidate order.

\texttt{m\allowbreak{}o\allowbreak{}s\allowbreak{}t\allowbreak{}V\allowbreak{}i\allowbreak{}v\allowbreak{}i\allowbreak{}d\allowbreak{}(\allowbreak{}c\allowbreak{}a\allowbreak{}n\allowbreak{}d\allowbreak{}i\allowbreak{}d\allowbreak{}a\allowbreak{}t\allowbreak{}e\allowbreak{}s\allowbreak{},\allowbreak{} \allowbreak{}o\allowbreak{}p\allowbreak{}t\allowbreak{}i\allowbreak{}o\allowbreak{}n\allowbreak{}s\allowbreak{})} chooses the highest-chroma candidate, with optional \texttt{a\allowbreak{}g\allowbreak{}a\allowbreak{}i\allowbreak{}n\allowbreak{}s\allowbreak{}t}, \texttt{m\allowbreak{}i\allowbreak{}n\allowbreak{}C\allowbreak{}o\allowbreak{}n\allowbreak{}t\allowbreak{}r\allowbreak{}a\allowbreak{}s\allowbreak{}t}, and \texttt{n\allowbreak{}o\allowbreak{}t} constraints. First-slice chroma is the deterministic RGB channel spread: \texttt{m\allowbreak{}a\allowbreak{}x\allowbreak{}(\allowbreak{}r\allowbreak{}e\allowbreak{}d\allowbreak{},\allowbreak{} \allowbreak{}g\allowbreak{}r\allowbreak{}e\allowbreak{}e\allowbreak{}n\allowbreak{},\allowbreak{} \allowbreak{}b\allowbreak{}l\allowbreak{}u\allowbreak{}e\allowbreak{})\allowbreak{} \allowbreak{}-\allowbreak{} \allowbreak{}m\allowbreak{}i\allowbreak{}n\allowbreak{}(\allowbreak{}r\allowbreak{}e\allowbreak{}d\allowbreak{},\allowbreak{} \allowbreak{}g\allowbreak{}r\allowbreak{}e\allowbreak{}e\allowbreak{}n\allowbreak{},\allowbreak{} \allowbreak{}b\allowbreak{}l\allowbreak{}u\allowbreak{}e\allowbreak{})}. If \texttt{m\allowbreak{}i\allowbreak{}n\allowbreak{}C\allowbreak{}o\allowbreak{}n\allowbreak{}t\allowbreak{}r\allowbreak{}a\allowbreak{}s\allowbreak{}t} is present, \texttt{a\allowbreak{}g\allowbreak{}a\allowbreak{}i\allowbreak{}n\allowbreak{}s\allowbreak{}t} is required.

\texttt{l\allowbreak{}e\allowbreak{}a\allowbreak{}s\allowbreak{}t\allowbreak{}V\allowbreak{}i\allowbreak{}v\allowbreak{}i\allowbreak{}d\allowbreak{}(\allowbreak{}c\allowbreak{}a\allowbreak{}n\allowbreak{}d\allowbreak{}i\allowbreak{}d\allowbreak{}a\allowbreak{}t\allowbreak{}e\allowbreak{}s\allowbreak{},\allowbreak{} \allowbreak{}o\allowbreak{}p\allowbreak{}t\allowbreak{}i\allowbreak{}o\allowbreak{}n\allowbreak{}s\allowbreak{})} chooses the lowest-chroma candidate, with the same optional constraints and the same \texttt{m\allowbreak{}i\allowbreak{}n\allowbreak{}C\allowbreak{}o\allowbreak{}n\allowbreak{}t\allowbreak{}r\allowbreak{}a\allowbreak{}s\allowbreak{}t} rule.

\texttt{c\allowbreak{}l\allowbreak{}o\allowbreak{}s\allowbreak{}e\allowbreak{}s\allowbreak{}t\allowbreak{}C\allowbreak{}o\allowbreak{}l\allowbreak{}o\allowbreak{}r\allowbreak{}(\allowbreak{}t\allowbreak{}a\allowbreak{}r\allowbreak{}g\allowbreak{}e\allowbreak{}t\allowbreak{},\allowbreak{} \allowbreak{}c\allowbreak{}a\allowbreak{}n\allowbreak{}d\allowbreak{}i\allowbreak{}d\allowbreak{}a\allowbreak{}t\allowbreak{}e\allowbreak{}s\allowbreak{})} chooses the closest candidate to a target color. The first slice uses squared RGB distance and keeps candidate order for equal distances.

\texttt{n\allowbreak{}t\allowbreak{}h\allowbreak{}C\allowbreak{}o\allowbreak{}l\allowbreak{}o\allowbreak{}r\allowbreak{}(\allowbreak{}c\allowbreak{}a\allowbreak{}n\allowbreak{}d\allowbreak{}i\allowbreak{}d\allowbreak{}a\allowbreak{}t\allowbreak{}e\allowbreak{}s\allowbreak{},\allowbreak{} \allowbreak{}i\allowbreak{}n\allowbreak{}d\allowbreak{}e\allowbreak{}x\allowbreak{})} chooses a candidate by ordered position. Integer indices use array-style ordering, negative integers count from the end, and out-of-range integers clamp to the nearest valid candidate. Fractional indices clamp to \texttt{0} through \texttt{1} and use \texttt{r\allowbreak{}o\allowbreak{}u\allowbreak{}n\allowbreak{}d\allowbreak{}(\allowbreak{}(\allowbreak{}l\allowbreak{}e\allowbreak{}n\allowbreak{}g\allowbreak{}t\allowbreak{}h\allowbreak{} \allowbreak{}-\allowbreak{} \allowbreak{}1\allowbreak{})\allowbreak{} \allowbreak{}*\allowbreak{} \allowbreak{}i\allowbreak{}n\allowbreak{}d\allowbreak{}e\allowbreak{}x\allowbreak{})}.

\section{Runtime Contract}\label{record-184-runtime-contract}

Selectors are pure. They do not read process state, filesystem data, terminal facts, or clocks.

Selector helpers return plain \texttt{T\allowbreak{}h\allowbreak{}e\allowbreak{}m\allowbreak{}e\allowbreak{}C\allowbreak{}o\allowbreak{}l\allowbreak{}o\allowbreak{}r\allowbreak{}R\allowbreak{}u\allowbreak{}l\allowbreak{}e\allowbreak{}D\allowbreak{}e\allowbreak{}f\allowbreak{}i\allowbreak{}n\allowbreak{}i\allowbreak{}t\allowbreak{}i\allowbreak{}o\allowbreak{}n} data. The token graph evaluates that intermediate rule definition, and the existing \texttt{T\allowbreak{}o\allowbreak{}k\allowbreak{}e\allowbreak{}n\allowbreak{}G\allowbreak{}r\allowbreak{}a\allowbreak{}p\allowbreak{}h\allowbreak{}.\allowbreak{}g\allowbreak{}e\allowbreak{}t\allowbreak{}(\allowbreak{})} path still returns a resolved \texttt{T\allowbreak{}o\allowbreak{}k\allowbreak{}e\allowbreak{}n\allowbreak{}V\allowbreak{}a\allowbreak{}l\allowbreak{}u\allowbreak{}e}.

Rule failures are deterministic:

\begin{itemize}
\tightlist
\item
  missing target, explicit candidate, candidate-scope, and exclusion paths throw with the missing path
\item
  vividness selectors throw when \texttt{m\allowbreak{}i\allowbreak{}n\allowbreak{}C\allowbreak{}o\allowbreak{}n\allowbreak{}t\allowbreak{}r\allowbreak{}a\allowbreak{}s\allowbreak{}t} is set without \texttt{a\allowbreak{}g\allowbreak{}a\allowbreak{}i\allowbreak{}n\allowbreak{}s\allowbreak{}t}
\item
  no minimum-contrast candidate throws unless the selector defines an explicit fallback
\item
  cycles remain cycle errors
\end{itemize}

Inspection dependencies include target and \texttt{a\allowbreak{}g\allowbreak{}a\allowbreak{}i\allowbreak{}n\allowbreak{}s\allowbreak{}t} references plus every candidate path read from explicit candidate lists or scopes, whether the candidate was selected, eligible-but-not-selected, excluded, invalid, or below the contrast floor.

\section{Accessibility And Assistive Posture}\label{record-184-accessibility-and-assistive-posture}

Selector rules improve color quality, but they do not make color sufficient as the only meaning carrier. The same graceful-lowering rule applies:

\begin{itemize}
\tightlist
\item
  no-color output must keep text, symbols, or structure
\item
  pipe output must not depend on swatches
\item
  accessible output must name state and intent
\end{itemize}

Contrast-aware selection is proof that a chosen color is readable when color is available. It is not proof that the UI is accessible by itself.

\section{Localization And Directionality Posture}\label{record-184-localization-and-directionality-posture}

No user-facing localized copy is required for the core rule primitives. Rule ids, token paths, candidate paths, and inspection facts are developer-facing identifiers.

Future DOGFOOD Theme Inspector prose that explains rule facts must flow through the DOGFOOD localization catalog. Token paths remain stable identifiers and are not localized.

\section{Agent Inspectability / Explainability Posture}\label{record-184-agent-inspectability-explainability-posture}

Inspection facts are a first-class product requirement. An agent should be able to assert:

\begin{itemize}
\tightlist
\item
  which selector rule resolved a token
\item
  which references and candidate scopes the rule read
\item
  which candidate won
\item
  which candidates were rejected
\item
  whether a candidate failed contrast, was excluded, or was the wrong type
\item
  which downstream token paths depend on the resolved token
\end{itemize}

These facts should not require terminal rendering, screenshot analysis, or RGB equality guessing.

\section{Linked Invariants}\label{record-184-linked-invariants}

\begin{itemize}
\tightlist
\item
  \href{https://github.com/flyingrobots/bijou/blob/cbd70612140e7a450d764b3c7be1c6af878b99bc/docs/invariants/runtime-truth-wins.md}{Runtime Truth Wins}: resolved values and inspection facts come from executing the graph.
\item
  \href{https://github.com/flyingrobots/bijou/blob/cbd70612140e7a450d764b3c7be1c6af878b99bc/docs/invariants/tests-are-the-spec.md}{Tests Are the Spec}: selector behavior lands through deterministic tests before docs claims.
\item
  \href{https://github.com/flyingrobots/bijou/blob/cbd70612140e7a450d764b3c7be1c6af878b99bc/docs/invariants/graceful-lowering-preserves-meaning.md}{Graceful Lowering Preserves Meaning}: color quality does not replace lower-mode structure.
\item
  \href{https://github.com/flyingrobots/bijou/blob/cbd70612140e7a450d764b3c7be1c6af878b99bc/docs/invariants/docs-are-the-demo.md}{Docs Are the Demo}: future Theme Lab surfaces should consume the same inspection facts tests assert.
\end{itemize}

\section{Implementation Outline}\label{record-184-implementation-outline}

\begin{enumerate}
\def\labelenumi{\arabic{enumi}.}
\tightlist
\item
  Add selector definition types to the core theme graph.
\item
  Collect deterministic candidates from explicit paths and scopes.
\item
  Implement contrast, vividness, closeness, and ordered selection helpers.
\item
  Expose \texttt{T\allowbreak{}o\allowbreak{}k\allowbreak{}e\allowbreak{}n\allowbreak{}G\allowbreak{}r\allowbreak{}a\allowbreak{}p\allowbreak{}h\allowbreak{}.\allowbreak{}i\allowbreak{}n\allowbreak{}s\allowbreak{}p\allowbreak{}e\allowbreak{}c\allowbreak{}t\allowbreak{}(\allowbreak{})} or an adjacent pure inspection API.
\item
  Preserve existing \texttt{T\allowbreak{}o\allowbreak{}k\allowbreak{}e\allowbreak{}n\allowbreak{}G\allowbreak{}r\allowbreak{}a\allowbreak{}p\allowbreak{}h\allowbreak{}.\allowbreak{}g\allowbreak{}e\allowbreak{}t\allowbreak{}(\allowbreak{})} behavior for raw, ref, mode, and transform definitions.
\item
  Wire the new types and helpers into the theme barrel and root package.
\item
  Document the recipe pattern in the design-system theme authoring docs.
\item
  Record the shipped behavior in \texttt{d\allowbreak{}o\allowbreak{}c\allowbreak{}s\allowbreak{}/\allowbreak{}C\allowbreak{}H\allowbreak{}A\allowbreak{}N\allowbreak{}G\allowbreak{}E\allowbreak{}L\allowbreak{}O\allowbreak{}G\allowbreak{}.\allowbreak{}m\allowbreak{}d}.
\end{enumerate}

\section{Tests To Write First}\label{record-184-tests-to-write-first}

\begin{itemize}
\tightlist
\item
  \texttt{b\allowbreak{}e\allowbreak{}s\allowbreak{}t\allowbreak{}C\allowbreak{}o\allowbreak{}n\allowbreak{}t\allowbreak{}r\allowbreak{}a\allowbreak{}s\allowbreak{}t\allowbreak{}W\allowbreak{}i\allowbreak{}t\allowbreak{}h\allowbreak{}(\allowbreak{})} chooses the readable foreground from a candidate scope.
\item
  \texttt{m\allowbreak{}i\allowbreak{}n\allowbreak{}C\allowbreak{}o\allowbreak{}n\allowbreak{}t\allowbreak{}r\allowbreak{}a\allowbreak{}s\allowbreak{}t\allowbreak{}W\allowbreak{}i\allowbreak{}t\allowbreak{}h\allowbreak{}(\allowbreak{})} rejects candidates below the minimum ratio and selects the first passing candidate.
\item
  \texttt{m\allowbreak{}o\allowbreak{}s\allowbreak{}t\allowbreak{}V\allowbreak{}i\allowbreak{}v\allowbreak{}i\allowbreak{}d\allowbreak{}(\allowbreak{})} respects \texttt{a\allowbreak{}g\allowbreak{}a\allowbreak{}i\allowbreak{}n\allowbreak{}s\allowbreak{}t}, \texttt{m\allowbreak{}i\allowbreak{}n\allowbreak{}C\allowbreak{}o\allowbreak{}n\allowbreak{}t\allowbreak{}r\allowbreak{}a\allowbreak{}s\allowbreak{}t}, and \texttt{n\allowbreak{}o\allowbreak{}t} exclusions.
\item
  \texttt{l\allowbreak{}e\allowbreak{}a\allowbreak{}s\allowbreak{}t\allowbreak{}V\allowbreak{}i\allowbreak{}v\allowbreak{}i\allowbreak{}d\allowbreak{}(\allowbreak{})} chooses the muted readable candidate.
\item
  \texttt{c\allowbreak{}l\allowbreak{}o\allowbreak{}s\allowbreak{}e\allowbreak{}s\allowbreak{}t\allowbreak{}C\allowbreak{}o\allowbreak{}l\allowbreak{}o\allowbreak{}r\allowbreak{}(\allowbreak{})} and \texttt{n\allowbreak{}t\allowbreak{}h\allowbreak{}(\allowbreak{})} are deterministic.
\item
  selector definitions participate in graph dependency inspection.
\item
  circular selector references still fail deterministically.
\item
  existing graph reference and transform tests keep passing.
\end{itemize}

\section{Acceptance Criteria}\label{record-184-acceptance-criteria}

\begin{itemize}
\tightlist
\item
  Bijou has no dependency on \texttt{d\allowbreak{}e\allowbreak{}s\allowbreak{}i\allowbreak{}g\allowbreak{}n\allowbreak{}-\allowbreak{}b\allowbreak{}o\allowbreak{}o\allowbreak{}k}.
\item
  Theme selectors are pure core functions.
\item
  Token graph rule definitions resolve to existing \texttt{T\allowbreak{}o\allowbreak{}k\allowbreak{}e\allowbreak{}n\allowbreak{}V\allowbreak{}a\allowbreak{}l\allowbreak{}u\allowbreak{}e} objects.
\item
  Inspection facts explain selected and rejected candidates.
\item
  Exclusions by token path prevent reserved roles from winning accent rules.
\item
  The new API exports from \texttt{@\allowbreak{}f\allowbreak{}l\allowbreak{}y\allowbreak{}i\allowbreak{}n\allowbreak{}g\allowbreak{}r\allowbreak{}o\allowbreak{}b\allowbreak{}o\allowbreak{}t\allowbreak{}s\allowbreak{}/\allowbreak{}b\allowbreak{}i\allowbreak{}j\allowbreak{}o\allowbreak{}u}.
\item
  Design-system docs show the primitive -\textgreater{} semantic rule -\textgreater{} UI token workflow.
\item
  Local validation includes focused theme tests, \texttt{n\allowbreak{}p\allowbreak{}m\allowbreak{} \allowbreak{}r\allowbreak{}u\allowbreak{}n\allowbreak{} \allowbreak{}t\allowbreak{}y\allowbreak{}p\allowbreak{}e\allowbreak{}c\allowbreak{}h\allowbreak{}e\allowbreak{}c\allowbreak{}k\allowbreak{}:\allowbreak{}t\allowbreak{}e\allowbreak{}s\allowbreak{}t}, \texttt{n\allowbreak{}p\allowbreak{}m\allowbreak{} \allowbreak{}r\allowbreak{}u\allowbreak{}n\allowbreak{} \allowbreak{}l\allowbreak{}i\allowbreak{}n\allowbreak{}t}, \texttt{n\allowbreak{}p\allowbreak{}m\allowbreak{} \allowbreak{}r\allowbreak{}u\allowbreak{}n\allowbreak{} \allowbreak{}c\allowbreak{}o\allowbreak{}d\allowbreak{}e\allowbreak{}-\allowbreak{}d\allowbreak{}o\allowbreak{}j\allowbreak{}o\allowbreak{}:\allowbreak{}c\allowbreak{}h\allowbreak{}a\allowbreak{}n\allowbreak{}g\allowbreak{}e\allowbreak{}d}, and \texttt{g\allowbreak{}i\allowbreak{}t\allowbreak{} \allowbreak{}d\allowbreak{}i\allowbreak{}f\allowbreak{}f\allowbreak{} \allowbreak{}-\allowbreak{}-\allowbreak{}c\allowbreak{}h\allowbreak{}e\allowbreak{}c\allowbreak{}k}.
\item
  Code Dojo debt remains verifiable through \texttt{n\allowbreak{}p\allowbreak{}m\allowbreak{} \allowbreak{}r\allowbreak{}u\allowbreak{}n\allowbreak{} \allowbreak{}c\allowbreak{}o\allowbreak{}d\allowbreak{}e\allowbreak{}-\allowbreak{}d\allowbreak{}o\allowbreak{}j\allowbreak{}o\allowbreak{}:\allowbreak{}d\allowbreak{}e\allowbreak{}b\allowbreak{}t} and \texttt{n\allowbreak{}p\allowbreak{}m\allowbreak{} \allowbreak{}r\allowbreak{}u\allowbreak{}n\allowbreak{} \allowbreak{}c\allowbreak{}o\allowbreak{}d\allowbreak{}e\allowbreak{}-\allowbreak{}d\allowbreak{}o\allowbreak{}j\allowbreak{}o\allowbreak{}:\allowbreak{}v\allowbreak{}e\allowbreak{}r\allowbreak{}i\allowbreak{}f\allowbreak{}y}; this cycle does not claim the repo goalpost unless that aggregate report reaches \texttt{1\allowbreak{}1\allowbreak{}2} or lower.
\end{itemize}

\section{Retrospective}\label{record-184-retrospective}

To be completed when the cycle lands.

